\chapter{Discussion générale}

\section{Ce que montre l’évaluation}

L’apport le plus solide d’Ariel Logminer ne réside pas dans le meilleur score prédictif. Il tient à la capacité de faire coexister des sources et des unités différentes tout en conservant des contrats et des preuves. Le routeur, le Contract Net, les magasins d’idempotence et les artefacts de campagne forment une chaîne inspectable. Cette propriété facilite l’analyse d’un résultat faible : il devient possible de déterminer si la difficulté vient du split, de la représentation, du modèle ou de l’orchestration.

Les évaluations prédictives confirment précisément la nécessité de cette prudence. Sur CICIDS2017, la différence entre split aléatoire et holdout scénario dépasse largement la variation entre modèles. Sur HDFS, le passage de l’événement au bloc transforme l’interprétation parce que le bloc est l’unité annotée. Sur BGL, le baseline de nouveauté révèle que le F1 global ne peut pas être lu sans la proportion de templates inconnus.

\section{Valeur de la spécialisation}

La spécialisation apporte d’abord une compatibilité technique. Un modèle tabulaire n’est pas appliqué directement à une séquence HDFS ; un parseur de templates n’est pas présenté comme classificateur de flux réseau. Le routeur automatise cette sélection et son rejet open-set évite une décision forcée pour les trois inconnus évalués.

Le bénéfice prédictif est plus nuancé. L’ablation du routage ne montre pas de gain systématique des spécialistes pour toutes les familles et graines. La conclusion défendable est donc que le routage fonctionne et protège la compatibilité des pipelines. Pour démontrer une amélioration prédictive, il faudrait un benchmark où modèle global et modèles spécialisés partagent exactement données, caractéristiques, budget d’optimisation et splits externes.

\section{Autonomie réelle mais bornée}

Les agents implémentés ne se limitent pas à consommer une file. Ils possèdent plusieurs handlers, annoncent des capacités, perçoivent leur état, calculent une utilité et peuvent refuser. Les messages \texttt{PROPOSE}, \texttt{REFUSE}, \texttt{AWARD} et \texttt{REJECT} rendent la décision visible. La campagne multi-VM observe quinze refus, un échec contrôlé et une réattribution réussie.

Cette autonomie reste déterministe et configurée. Les coefficients de l’utilité ne sont pas appris. Le coordinateur conserve la décision finale d’attribution. Parler d’un système pleinement auto-adaptatif dépasserait donc les preuves. Le terme « autonome » garde un sens précis : décision locale bornée dans un protocole explicite.

\section{Coût de l’orchestration}

Le benchmark A--D montre que les agents ne sont pas une optimisation gratuite. Sur les micro-tâches, le monolithe évite les messages, la sérialisation et la sélection des offres. Les architectures autonomes paient ce coût et la mémoire activée ne produit pas d’amélioration systématique.

Ce résultat n’invalide pas l’architecture. Il distingue deux objectifs souvent confondus : minimiser le temps d’une opération élémentaire et rendre une décision distribuée observable. Ariel Logminer privilégie la modularité, le refus explicite, la reprise et la traçabilité. Une évaluation future sur des tâches d’inférence plus longues devrait mesurer si le coût fixe devient secondaire, mais aucune conclusion anticipée n’est ajoutée au présent mémoire.

\section{Mémoire et apprentissage}

La mémoire influence l’utilité de l’agent dans des scénarios contrôlés. Elle constitue donc un mécanisme fonctionnel d’adaptation comportementale. Le benchmark ne montre toutefois aucun gain systématique de débit, de latence ou d’équité. Son intérêt actuel est d’offrir un état auditable pour des politiques futures, non d’améliorer automatiquement les métriques.

La gouvernance des modèles appelle la même distinction. Le dépôt contient le code de comparaison et le dry-run ; certaines branches sont testées. Une campagne complète de promotion et rejet, suivie d’une mesure avant/après sur jeu gelé, n’est pas documentée. La mise à jour contrôlée est partiellement implémentée, mais son bénéfice reste non évalué.

\section{Portée opérationnelle}

La couverture multi-format et la campagne M montrent que la chaîne peut lire, normaliser et traiter plusieurs sources dans le laboratoire. Le dashboard donne accès aux états, alertes et incidents. Ces éléments forment un prototype utilisable pour l’expérimentation et la démonstration.

Ils ne constituent pas un produit SOC. La topologie ne possède ni redondance Redis démontrée, ni contrôle d’accès complet, ni gestion industrielle des secrets, ni observabilité distribuée à longue durée. Les performances locales ne sont pas des objectifs de service. Le système doit être présenté comme une base d’ingénierie dont les frontières sont connues.

\section{Réponses transversales}

Trois enseignements se dégagent. Premièrement, la traçabilité du protocole est une condition de la qualité scientifique : un score sans split ou unité est insuffisant. Deuxièmement, l’architecture multi-agent apporte une capacité d’attribution et de reprise, mais pas un gain de performance automatique. Troisièmement, les contrôles négatifs sont informatifs : la faiblesse du holdout CICIDS2017, l’absence de positifs BGL connus et l’instabilité sous permutation empêchent une présentation trop favorable.

Le mémoire assume donc une contribution d’intégration. Il montre comment assembler et éprouver des briques connues en conservant la possibilité de réfuter une affirmation. Cette orientation est cohérente avec un travail d’ingénierie : la crédibilité du système vient autant de ses limites documentées que de ses chemins réussis.

\section{Priorités de prolongement}

Les travaux les plus utiles seraient, dans l’ordre, d’élargir les tests open-set, d’exécuter des validations temporelles sur davantage de périodes, de constituer un corpus Apache réel et de mesurer le Contract Net sur des tâches représentatives. Une campagne de gouvernance devrait ensuite provoquer des candidats meilleurs et moins bons, observer promotion et rejet, puis mesurer l’effet sur un test gelé.

L’industrialisation viendrait après ces contrôles scientifiques : réplication et basculement Redis, partitions réseau, sécurité des échanges, secrets, métriques opérationnelles et tests prolongés. Ces éléments sont des perspectives et non des propriétés actuelles d’Ariel Logminer.
